\chapter{El movimiento: uniforme y variado}\label{ch:g9:uniform-varied-motion}

Hace dos años diagnosticabas el movimiento a ojo: puntos igual de
separados, uniforme; puntos que se separaban, acelerando. Este año los
puntos reciben rayas de milímetro y el diagnóstico recibe números ---
intervalo a intervalo, \omterm{def:g7:motion-average-speed:formula}{velocidad} a \omterm{def:g7:motion-average-speed:formula}{velocidad}. Y al final
del capítulo aparece una frase tranquila sobre la
\omterm{def:g2:pushes-and-pulls:force}{fuerza} y el movimiento que algún día llevará media física a
cuestas.

\section{El registro del movimiento}

\begin{definition}[Registros de posiciones cronometradas]\label{def:g9:uniform-varied-motion:record}
Un \emph{registro de posiciones cronometradas}\index{registro de
posiciones cronometradas} de un movimiento marca las posiciones del
objeto que se mueve a intervalos iguales de tiempo --- fotogramas de
vídeo pasados uno a uno, una fotografía con destellos estroboscópicos o
el marcador de puntos del laboratorio que estampa cincuenta puntos por
segundo en una cinta de papel tirada por el objeto. Entre dos marcas
vecinas cualesquiera está el viaje de un intervalo: el registro
convierte el movimiento en tramos medibles.
\end{definition}

\begin{method}[Leer un registro con números]\label{met:g9:uniform-varied-motion:read}
Para un registro con intervalos de \omterm{def:g2:measuring-time:duration}{duración} $\tau$ (por ejemplo
\qty{0.02}{s} para cincuenta por segundo, o unos cómodos \qty{0.1}{s}
de vídeo):
\begin{enumerate}
  \item mide con la regla milimetrada cada hueco entre marcas
        vecinas;
  \item divide cada hueco entre el intervalo: $v = d/\tau$ da la
        \omterm{def:g7:motion-average-speed:formula}{velocidad} (media) en ese pequeño tramo --- en \unit{m/s}
        si los huecos van en \omterm{def:g2:measuring-length:units}{metros};
  \item pon las \omterm{def:g7:motion-average-speed:formula}{velocidades} una al lado de otra y lee la
        historia: números constantes, \omterm{def:g6:describing-motion:uniform}{movimiento uniforme}; que
        suben, acelerado; que bajan, decelerado;
  \item cita la \omterm{def:g7:motion-average-speed:formula}{velocidad} de cualquier tramo como la
        \omterm{def:g7:motion-average-speed:formula}{velocidad} ``en'' ese instante: en un intervalo tan
        corto, la media y la instantánea se dan la mano.
\end{enumerate}
\end{method}

\begin{omfigure}
\begin{tikzpicture}[scale=0.85]
  \foreach \x in {0.5,2.0,3.5,5.0,6.5,8.0,9.5} {
    \fill[omDef] (\x,2.6) circle (0.1);
  }
  \node[right, font=\small] at (10.0,2.6)
    {huecos todos de \qty{3.0}{cm}: uniforme};
  \foreach \x in {0.5,1.1,2.0,3.2,4.7,6.5,8.6} {
    \fill[omThm] (\x,1.5) circle (0.1);
  }
  \node[right, font=\small] at (10.0,1.5)
    {huecos que crecen: acelerado};
  \foreach \x in {0.5,2.6,4.3,5.6,6.5,7.0,7.2} {
    \fill[omMeth] (\x,0.4) circle (0.1);
  }
  \node[right, font=\small] at (10.0,0.4)
    {huecos que encogen: decelerado};
\end{tikzpicture}

{\small Tres cintas del marcador de puntos, un punto por intervalo.
Este año los huecos se miden, y cada uno se convierte en una
\omterm{def:g7:motion-average-speed:formula}{velocidad}.}
\end{omfigure}

\begin{example}[Una cinta, resuelta]\label{ex:g9:uniform-varied-motion:worked}
La cinta de un carrito, con intervalos de \qty{0.1}{s} y huecos en
\omterm{def:g2:measuring-length:units}{centímetros}: $2.0$, $3.0$, $4.0$, $5.0$, $6.0$.
\omterm{def:g7:motion-average-speed:formula}{Velocidades}, hueco a hueco: $0.02 \div 0.1 = \qty{0.2}{m/s}$, y
después $0.3$, $0.4$, $0.5$ y \qty{0.6}{m/s}. Veredicto: acelerado ---
y con una regularidad preciosa: la \omterm{def:g7:motion-average-speed:formula}{velocidad} crece exactamente
\qty{0.1}{m/s} en cada intervalo. Los movimientos cuya
\omterm{def:g7:motion-average-speed:formula}{velocidad} sube a pasos iguales en tiempos iguales tienen un
nombre especial y un gran futuro --- \emph{uniformemente} acelerados,
la firma de las causas constantes.
\end{example}

\begin{example}[El movimiento acelerado más famoso]\label{ex:g9:uniform-varied-motion:freefall}
Suelta una piedra y grábala: los huecos se estiran intervalo a
intervalo --- la caída libre es acelerada. Medida, su
\omterm{def:g7:motion-average-speed:formula}{velocidad} crece unos \qty{9.8}{m/s} en cada segundo de caída:
después de un segundo, \qty{9.8}{m/s}; después de dos, casi
\qty{20}{m/s}. Que ese ritmo de crecimiento sea igual al $g$ del lugar
no es casualidad --- es la rima más profunda de la física de este año,
y el volumen de bachillerato construye su mecánica sobre ella. Anótala,
subráyala y déjala esperar.
\end{example}

\section{La velocidad instantánea, con honradez}

\begin{definition}[Velocidad instantánea]\label{def:g9:uniform-varied-motion:instant}
La \emph{velocidad instantánea}\index{velocidad instantánea} de un
movimiento en un instante es la \omterm{def:g7:motion-average-speed:average}{velocidad media} en un intervalo
muy corto alrededor de ese instante --- tan corto que el movimiento no
tiene sitio para cambiar de ritmo dentro de él. Es el número del
velocímetro y el del radar: los dos, de hecho, miden exactamente
como el \cref{met:g9:uniform-varied-motion:read}, con intervalos de un
parpadeo. (Qué significa ``muy corto'' con todo el rigor es la gran
pregunta que responden las matemáticas del último año de bachillerato
--- y la respuesta creó la ciencia moderna.)
\end{definition}

\begin{example}[La media y la instantánea, lado a lado]\label{ex:g9:uniform-varied-motion:sidebyside}
Un trayecto de metro entre estaciones: \omterm{def:g7:motion-average-speed:average}{velocidad media} ---
distancia total entre tiempo total --- \qty{35}{km/h}.
\omterm{def:g9:uniform-varied-motion:instant}{Velocidad instantánea}: $0$ en los dos andenes, \qty{70}{km/h} a
mitad de túnel. En el \omterm{def:g6:describing-motion:uniform}{movimiento uniforme}, y solo ahí, las dos
nociones se funden en un único número constante: el
\omterm{def:g6:describing-motion:uniform}{movimiento uniforme} es precisamente el movimiento que convierte la
\omterm{def:g7:motion-average-speed:formula}{velocidad} en una magnitud única, honrada y sin intervalos.
\end{example}

\section{La frase tranquila}

\begin{proposition}[Las fuerzas equilibradas dejan el movimiento como está]\label{prop:g9:uniform-varied-motion:inertia}
Cuando las \omterm{def:g2:pushes-and-pulls:force}{fuerzas} sobre un cuerpo se cancelan entre sí --- o
cuando no actúa ninguna ---, el movimiento del cuerpo no cambia: en
reposo sigue en reposo, y en movimiento continúa en línea recta a
\omterm{def:g7:motion-average-speed:formula}{velocidad} constante. El cambio de movimiento --- acelerar,
frenar, girar --- solo ocurre mientras actúa alguna
\omterm{def:g2:pushes-and-pulls:force}{fuerza} \emph{no \omterm{def:g2:balance-equilibrium:equilibrium}{equilibrada}}. Esta frase, enunciada aquí y
honradamente aplazada, es la ley inaugural de toda la mecánica; el
volumen de bachillerato empieza con ella y ya no la suelta.
\end{proposition}

\begin{example}[Leer el mundo con la frase]\label{ex:g9:uniform-varied-motion:reading}
El disco sobre el hielo liso se desliza recto y constante ---
\omterm{def:g2:pushes-and-pulls:force}{fuerzas} \omterm{def:g2:balance-equilibrium:equilibrium}{equilibradas} (\omterm{def:g4:weight-and-mass:weight}{peso} hacia abajo, hielo hacia
arriba), movimiento sin cambios: no hace falta ningún empujón para
\emph{seguir} yendo. El coche que frena aminora: una
\omterm{def:g2:pushes-and-pulls:force}{fuerza} no \omterm{def:g2:balance-equilibrium:equilibrium}{equilibrada}, el agarre de la carretera sobre los
neumáticos, actuando hacia atrás. El ciclista que gira:
\omterm{def:g2:pushes-and-pulls:force}{fuerza} no \omterm{def:g2:balance-equilibrium:equilibrium}{equilibrada} hacia el lado, aportada por los neumáticos
inclinados. Y el avión de línea en crucero a \qty{900}{km/h} constantes:
el empuje equilibra la \omterm{def:g8:ohms-law:resistance}{resistencia}, la sustentación equilibra el
\omterm{def:g4:weight-and-mass:weight}{peso} --- cuatro \omterm{def:g2:pushes-and-pulls:force}{fuerzas}, empate perfecto,
\omterm{def:g6:describing-motion:uniform}{movimiento uniforme}. La quietud y el crucero son, para la física,
el mismo estado apacible.
\end{example}

\begin{remark}[Por qué sorprende la frase]\label{rem:g9:uniform-varied-motion:surprise}
La vida diaria parece protestar: deja de pedalear y la bicicleta
aminora --- ¿no \emph{necesita} entonces \omterm{def:g2:pushes-and-pulls:force}{fuerza} el movimiento?
Pero la bicicleta que aminora no está libre de \omterm{def:g2:pushes-and-pulls:force}{fuerzas}: el rozamiento y
el \omterm{def:g2:air-around-us:air}{aire} empujan hacia atrás, sin equilibrar, y por eso exactamente
decae el ritmo. Quítalos --- la pista de hielo, el vacío del espacio ---
y el movimiento se desliza para siempre: las sondas espaciales lanzadas
antes de que se conocieran tus abuelos siguen deslizándose ahora, con
los \omterm{ex:g6:circuits-and-safety:motor}{motores} fríos. La humanidad necesitó dos mil años para ver a través
del disfraz del rozamiento; tú tienes la ventaja de las pistas de hielo
y de las sondas espaciales.
\end{remark}

\section{Ejercicios}

\begin{exercise}[$\star$]\label{exo:g9:uniform-varied-motion:1}
¿Qué es un \omterm{def:g9:uniform-varied-motion:record}{registro de posiciones cronometradas}? Nombra dos maneras
de hacer uno.
\end{exercise}

\begin{exercise}[$\star$]\label{exo:g9:uniform-varied-motion:2}
Una cinta con intervalos de \qty{0.1}{s} enseña huecos de
\qty{4.0}{cm}, $4.0$, $4.0$ y $4.0$. Diagnostica el movimiento y da su
\omterm{def:g7:motion-average-speed:formula}{velocidad} en \unit{m/s}.
\end{exercise}

\begin{exercise}[$\star$]\label{exo:g9:uniform-varied-motion:3}
Otra cinta: huecos de $1.0$, $2.5$, $4.5$ y \qty{7.0}{cm}. Diagnostica
--- y calcula las \omterm{def:g7:motion-average-speed:formula}{velocidades} del primer tramo y del último.
\end{exercise}

\begin{exercise}[$\star$]\label{exo:g9:uniform-varied-motion:4}
Distingue la media de la \omterm{def:g9:uniform-varied-motion:instant}{velocidad instantánea} en el camino al
colegio: ¿cuál lee el radar de la entrada y cuál calcula el ``veinte
minutos de puerta a puerta'' de la familia?
\end{exercise}

\begin{exercise}[$\star$]\label{exo:g9:uniform-varied-motion:5}
¿En qué movimientos coinciden la media y la \omterm{def:g9:uniform-varied-motion:instant}{velocidad
instantánea}? ¿Por qué ahí y solo ahí?
\end{exercise}

\begin{exercise}[$\star$]\label{exo:g9:uniform-varied-motion:6}
¿Cuánto crece cada segundo la \omterm{def:g7:motion-average-speed:formula}{velocidad} de una piedra en caída
libre, y a qué número local es igual ese crecimiento?
\end{exercise}

\begin{exercise}[$\star$]\label{exo:g9:uniform-varied-motion:7}
Enuncia la frase tranquila. ¿Qué dice sobre un cuerpo sin ninguna
\omterm{def:g2:pushes-and-pulls:force}{fuerza} --- y sobre el empate a cuatro del avión en crucero?
\end{exercise}

\begin{exercise}[$\star$]\label{exo:g9:uniform-varied-motion:8}
``El movimiento necesita \omterm{def:g2:pushes-and-pulls:force}{fuerza} para continuar''. Condena este
error antiguo con la sonda espacial que se desliza y con el disfraz del
rozamiento.
\end{exercise}

\begin{exercise}[$\star\star$]\label{exo:g9:uniform-varied-motion:9}
Una cinta (intervalos de \qty{0.1}{s}) marca, en
\omterm{def:g2:measuring-length:units}{centímetros}: $6.0$, $5.0$, $4.0$, $3.0$, $2.0$. Diagnostica;
calcula la \omterm{def:g7:motion-average-speed:formula}{velocidad} de cada tramo; y predice el futuro de la
cinta si la pauta se mantiene.
\end{exercise}

\begin{exercise}[$\star\star$]\label{exo:g9:uniform-varied-motion:10}
El carrito resuelto ganaba \qty{0.1}{m/s} por intervalo de
\qty{0.1}{s}. Expresa ese crecimiento por \emph{segundo} --- y compara
la ganancia del carrito con la de la caída libre.
\end{exercise}

\begin{exercise}[$\star\star$]\label{exo:g9:uniform-varied-motion:11}
Una paracaidista cae cada vez más rápido --- y después, con el
paracaídas abierto, a \qty{5}{m/s} constantes. Lee las dos
\omterm{def:g7:moon-phases:phases}{fases} con la frase tranquila: ¿qué está sin equilibrar al
principio y qué empate se ha formado al final?
\end{exercise}

\begin{exercise}[$\star\star\star$]\label{exo:g9:uniform-varied-motion:12}
Diseña el veredicto completo de laboratorio sobre un cochecito de
juguete soltado por una rampa y lanzado sobre una alfombra: el registro
que hay que tomar, los números que hay que calcular, el diagnóstico
previsto en dos actos (qué acto en la rampa y cuál en la alfombra) y la
lectura de cada acto en el lenguaje de la frase de las \omterm{def:g2:pushes-and-pulls:force}{fuerzas}.
\end{exercise}

\section{Problema: la oficina de las cintas}

\begin{problem}[{Problema de fin de semana --- una tarde en la oficina
de análisis del movimiento; cuatro cintas, cuatro veredictos; el
ripio del inspector}]\label{pb:g9:uniform-varied-motion:1}
La oficina analiza registros cronometrados para colegios, clubes
deportivos y juzgados. Los intervalos son de \qty{0.1}{s} en todo el
problema y los huecos llegan en \omterm{def:g2:measuring-length:units}{centímetros}. Tú manejas la
regla.

\medskip\noindent\textbf{Parte I --- Veredictos de rutina.}
\begin{enumerate}
  \item Cinta A (alguien que camina por un pasillo): $8.0$, $8.0$,
        $8.0$, $8.0$, $8.0$. Veredicto y \omterm{def:g7:motion-average-speed:formula}{velocidad} --- en
        \unit{m/s} y en \unit{km/h}.
  \item Cinta B (la salida de un velocista): $2.0$, $4.0$, $6.0$,
        $8.0$, $10.0$. Veredicto, \omterm{def:g7:motion-average-speed:formula}{velocidades} primera y última,
        y el crecimiento de \omterm{def:g7:motion-average-speed:formula}{velocidad} por segundo.
  \item Cinta C (una bola que rueda y se encuentra con arena):
        $10.0$, $7.0$, $4.5$, $2.5$, $1.0$. Veredicto --- y la
        explicación de la arena según la frase de las \omterm{def:g2:pushes-and-pulls:force}{fuerzas}.
  \item Cinta D (un misterio): $5.0$, $5.0$, $5.0$, $7.5$, $10.0$.
        Cuenta la historia en dos actos de D y marca el intervalo en
        el que pasó algo.
\end{enumerate}

\medskip\noindent\textbf{Parte II --- La mesa de física de la
oficina.}
\begin{enumerate}[resume]
  \item Un cliente insiste en que su moto de reparto ``iba a
        \qty{30}{km/h}, agente, de media''. El radar registró
        \qty{55}{km/h} en el cruce. Explícale al cliente por qué los
        dos números pueden ser ciertos y cuál es el que lee la ley.
  \item El velocista de la cinta B gana \omterm{def:g7:motion-average-speed:formula}{velocidad} a pasos
        iguales. ¿Qué palabra única estampa la oficina sobre un
        movimiento así, y qué sugiere la regularidad del crecimiento
        sobre la \omterm{def:g2:pushes-and-pulls:force}{fuerza} que empuja?
  \item Quien camina por el pasillo de la cinta A, dice la frase de
        las \omterm{def:g2:pushes-and-pulls:force}{fuerzas}, camina con las \omterm{def:g2:pushes-and-pulls:force}{fuerzas} \omterm{def:g2:balance-equilibrium:equilibrium}{equilibradas}. Nombra
        el \omterm{def:g2:balance-equilibrium:equilibrium}{equilibrio} de un caminante a ritmo constante (qué empuja y
        qué se resiste).
  \item La cinta de caída de una piedra enseña \omterm{def:g7:motion-average-speed:formula}{velocidades} de
        tramo de $0.98$, $1.96$ y \qty{2.94}{m/s}. Verifica el
        crecimiento por segundo y nombra el número que reproduce.
\end{enumerate}

\medskip\noindent\textbf{Parte III --- El caso judicial.}
Un monopatín rodó desde una rampa por el patio del colegio hasta un
arriate; el conserje culpa a la ``\omterm{def:g7:motion-average-speed:formula}{velocidad} temeraria'' y el
patinador afirma que ``iba frenando todo el rato''. La cámara del patio
da estos huecos: salida de la rampa $12.0$ y después $11.5$, $11.0$,
$10.5$ y $10.0$ al entrar en el arriate.
\begin{enumerate}[resume]
  \item Veredicto sobre la afirmación del patinador --- con las
        \omterm{def:g7:motion-average-speed:formula}{velocidades} de cada tramo.
  \item \omterm{def:g7:motion-average-speed:formula}{Velocidad} de salida y \omterm{def:g7:motion-average-speed:formula}{velocidad} de entrada en el
        arriate, en \unit{km/h}: ¿alguna de las dos era ``temeraria''
        al lado de los \qty{15}{km/h} de un ciclista con brío?
  \item El conserje pregunta por qué el monopatín no se paró sin más
        en el suelo llano, ``como se paran las cosas de forma
        natural''. La respuesta de la oficina, en una frase de
        lenguaje de \omterm{def:g2:pushes-and-pulls:force}{fuerzas}.
  \item Cierra el expediente con el ripio del inspector --- dos
        líneas: qué significan los huecos iguales y qué exigen los
        huecos que cambian (una causa, según la frase tranquila).
\end{enumerate}
\end{problem}
